% Options for packages loaded elsewhere
\PassOptionsToPackage{unicode}{hyperref}
\PassOptionsToPackage{hyphens}{url}
\PassOptionsToPackage{dvipsnames,svgnames,x11names}{xcolor}
\PassOptionsToPackage{space}{xeCJK}
%
\documentclass[
  11pt,
  a4paper,
]{article}
\usepackage{amsmath,amssymb}
\usepackage{iftex}
\ifPDFTeX
  \usepackage[T1]{fontenc}
  \usepackage[utf8]{inputenc}
  \usepackage{textcomp} % provide euro and other symbols
\else % if luatex or xetex
  \usepackage{unicode-math} % this also loads fontspec
  \defaultfontfeatures{Scale=MatchLowercase}
  \defaultfontfeatures[\rmfamily]{Ligatures=TeX,Scale=1}
\fi
\usepackage{lmodern}
\ifPDFTeX\else
  % xetex/luatex font selection
  \setmainfont[]{TeX Gyre Termes}
  \setmonofont[]{DejaVu Sans Mono}
  \ifXeTeX
    \usepackage{xeCJK}
    \setCJKmainfont[]{Noto Serif CJK SC}
          \fi
  \ifLuaTeX
    \usepackage[]{luatexja-fontspec}
    \setmainjfont[]{Noto Serif CJK SC}
  \fi
\fi
% Use upquote if available, for straight quotes in verbatim environments
\IfFileExists{upquote.sty}{\usepackage{upquote}}{}
\IfFileExists{microtype.sty}{% use microtype if available
  \usepackage[]{microtype}
  \UseMicrotypeSet[protrusion]{basicmath} % disable protrusion for tt fonts
}{}
\makeatletter
\@ifundefined{KOMAClassName}{% if non-KOMA class
  \IfFileExists{parskip.sty}{%
    \usepackage{parskip}
  }{% else
    \setlength{\parindent}{0pt}
    \setlength{\parskip}{6pt plus 2pt minus 1pt}}
}{% if KOMA class
  \KOMAoptions{parskip=half}}
\makeatother
\usepackage{xcolor}
\usepackage[margin=22mm]{geometry}
\usepackage{color}
\usepackage{fancyvrb}
\newcommand{\VerbBar}{|}
\newcommand{\VERB}{\Verb[commandchars=\\\{\}]}
\DefineVerbatimEnvironment{Highlighting}{Verbatim}{commandchars=\\\{\}}
% Add ',fontsize=\small' for more characters per line
\newenvironment{Shaded}{}{}
\newcommand{\AlertTok}[1]{\textcolor[rgb]{1.00,0.00,0.00}{\textbf{#1}}}
\newcommand{\AnnotationTok}[1]{\textcolor[rgb]{0.38,0.63,0.69}{\textbf{\textit{#1}}}}
\newcommand{\AttributeTok}[1]{\textcolor[rgb]{0.49,0.56,0.16}{#1}}
\newcommand{\BaseNTok}[1]{\textcolor[rgb]{0.25,0.63,0.44}{#1}}
\newcommand{\BuiltInTok}[1]{\textcolor[rgb]{0.00,0.50,0.00}{#1}}
\newcommand{\CharTok}[1]{\textcolor[rgb]{0.25,0.44,0.63}{#1}}
\newcommand{\CommentTok}[1]{\textcolor[rgb]{0.38,0.63,0.69}{\textit{#1}}}
\newcommand{\CommentVarTok}[1]{\textcolor[rgb]{0.38,0.63,0.69}{\textbf{\textit{#1}}}}
\newcommand{\ConstantTok}[1]{\textcolor[rgb]{0.53,0.00,0.00}{#1}}
\newcommand{\ControlFlowTok}[1]{\textcolor[rgb]{0.00,0.44,0.13}{\textbf{#1}}}
\newcommand{\DataTypeTok}[1]{\textcolor[rgb]{0.56,0.13,0.00}{#1}}
\newcommand{\DecValTok}[1]{\textcolor[rgb]{0.25,0.63,0.44}{#1}}
\newcommand{\DocumentationTok}[1]{\textcolor[rgb]{0.73,0.13,0.13}{\textit{#1}}}
\newcommand{\ErrorTok}[1]{\textcolor[rgb]{1.00,0.00,0.00}{\textbf{#1}}}
\newcommand{\ExtensionTok}[1]{#1}
\newcommand{\FloatTok}[1]{\textcolor[rgb]{0.25,0.63,0.44}{#1}}
\newcommand{\FunctionTok}[1]{\textcolor[rgb]{0.02,0.16,0.49}{#1}}
\newcommand{\ImportTok}[1]{\textcolor[rgb]{0.00,0.50,0.00}{\textbf{#1}}}
\newcommand{\InformationTok}[1]{\textcolor[rgb]{0.38,0.63,0.69}{\textbf{\textit{#1}}}}
\newcommand{\KeywordTok}[1]{\textcolor[rgb]{0.00,0.44,0.13}{\textbf{#1}}}
\newcommand{\NormalTok}[1]{#1}
\newcommand{\OperatorTok}[1]{\textcolor[rgb]{0.40,0.40,0.40}{#1}}
\newcommand{\OtherTok}[1]{\textcolor[rgb]{0.00,0.44,0.13}{#1}}
\newcommand{\PreprocessorTok}[1]{\textcolor[rgb]{0.74,0.48,0.00}{#1}}
\newcommand{\RegionMarkerTok}[1]{#1}
\newcommand{\SpecialCharTok}[1]{\textcolor[rgb]{0.25,0.44,0.63}{#1}}
\newcommand{\SpecialStringTok}[1]{\textcolor[rgb]{0.73,0.40,0.53}{#1}}
\newcommand{\StringTok}[1]{\textcolor[rgb]{0.25,0.44,0.63}{#1}}
\newcommand{\VariableTok}[1]{\textcolor[rgb]{0.10,0.09,0.49}{#1}}
\newcommand{\VerbatimStringTok}[1]{\textcolor[rgb]{0.25,0.44,0.63}{#1}}
\newcommand{\WarningTok}[1]{\textcolor[rgb]{0.38,0.63,0.69}{\textbf{\textit{#1}}}}
\usepackage{longtable,booktabs,array}
\usepackage{calc} % for calculating minipage widths
% Correct order of tables after \paragraph or \subparagraph
\usepackage{etoolbox}
\makeatletter
\patchcmd\longtable{\par}{\if@noskipsec\mbox{}\fi\par}{}{}
\makeatother
% Allow footnotes in longtable head/foot
\IfFileExists{footnotehyper.sty}{\usepackage{footnotehyper}}{\usepackage{footnote}}
\makesavenoteenv{longtable}
\setlength{\emergencystretch}{3em} % prevent overfull lines
\providecommand{\tightlist}{%
  \setlength{\itemsep}{0pt}\setlength{\parskip}{0pt}}
\setcounter{secnumdepth}{-\maxdimen} % remove section numbering
\ifLuaTeX
\usepackage[bidi=basic]{babel}
\else
\usepackage[bidi=default]{babel}
\fi
\babelprovide[main,import]{chinese}
\ifPDFTeX
\else
\babelfont{rm}[]{TeX Gyre Termes}
\fi
% get rid of language-specific shorthands (see #6817):
\let\LanguageShortHands\languageshorthands
\def\languageshorthands#1{}
\setCJKsansfont{Noto Sans CJK SC}
\setCJKmonofont{Noto Sans CJK SC}
\ifLuaTeX
  \usepackage{selnolig}  % disable illegal ligatures
\fi
\IfFileExists{bookmark.sty}{\usepackage{bookmark}}{\usepackage{hyperref}}
\IfFileExists{xurl.sty}{\usepackage{xurl}}{} % add URL line breaks if available
\urlstyle{same}
\hypersetup{
  pdftitle={仅迭代一次项的变分重构与 NCFV 高效通量耦合},
  pdfauthor={DNDSR 数学与实现设计说明},
  pdflang={zh-CN},
  colorlinks=true,
  linkcolor={Maroon},
  filecolor={Maroon},
  citecolor={Blue},
  urlcolor={Blue},
  pdfcreator={LaTeX via pandoc}}

\title{仅迭代一次项的变分重构与 NCFV 高效通量耦合}
\usepackage{etoolbox}
\makeatletter
\providecommand{\subtitle}[1]{% add subtitle to \maketitle
  \apptocmd{\@title}{\par {\large #1 \par}}{}{}
}
\makeatother
\subtitle{严格 Schur 消元、梯度闭合降维及可实施的并行流程}
\author{DNDSR 数学与实现设计说明}
\date{2026 年 9 月 28 日}

\begin{document}
\maketitle

\hypertarget{ux53efux4ee5ux53eaux8fedux4ee3ux4e00ux6b21ux9879ux4f46ux9700ux8981ux4fddux7559ux4e8cux6b21ux9879ux7684ux6570ux5b66ux4f5cux7528}{%
\section{1.
可以只迭代一次项，但需要保留二次项的数学作用}\label{ux53efux4ee5ux53eaux8fedux4ee3ux4e00ux6b21ux9879ux4f46ux9700ux8981ux4fddux7559ux4e8cux6b21ux9879ux7684ux6570ux5b66ux4f5cux7528}}

\textbf{可以把一次项作为唯一的外层未知量。要做到这一点，必须消去二次项，或用一次项给二次项建立明确的闭合关系。直接把二次列删掉、把二次系数设为零，一般不能保持二次多项式精确性。}

本文给出两条完整路线。路线 A 对原变分方程作全局 Schur
消元，得到与原方程严格等价的一次项系统；它仍需处理二次项的全局逆作用。路线
B 用邻点梯度差恢复
Hessian，再将该关系代入原泛函后重新求变分，得到只有一次项未知量的稀疏系统；它保留二次场再现能力，但对一般数据属于一个新的受约束变分重构。

二者都可以接入 NCFV
高效通量积分，因为高效积分只需要节点点值和一阶梯度。高效积分不用
Hessian，并不自动意味着传统变分重构迭代也可以省略 Hessian 的影响。

截至本说明编写时，当前源码 \texttt{ComputeCoefficients} 仍迭代三维每节点
9 个二次系数，迭代结束后提取前 3
项并除以参考长度，把物理梯度交给高效通量。\texttt{RecoverPointValues}
在高效模式下已用梯度积分权重恢复点值。此前的界面面均值变分/PCG
实现已在工作区标记为历史方案；本说明以当前边中点变分算子和高效积分接口为基础。本次交付为推导及实现设计，尚未把下面两种降维方法写入求解器，也没有它们的新精度数据。

\hypertarget{ux4e00ux6b21ux9879ux4e8cux6b21ux9879ux4e0eux7269ux7406ux5bfcux6570}{%
\section{2.
一次项、二次项与物理导数}\label{ux4e00ux6b21ux9879ux4e8cux6b21ux9879ux4e0eux7269ux7406ux5bfcux6570}}

考虑空间维数 \(d\)。已知节点对偶体积均值 \(\bar u_i\)，体积
\(V_i\)，节点位置 \(x_i\)，分方向尺度矩阵
\(L_i=\operatorname{diag}(L_{i1},\ldots,L_{id})\)。令
\(\xi_i=L_i^{-1}(x-x_i)\)，取零均值线性和二次基：

\[
P_i(x)=\bar u_i+\phi_i^{(1)}(x)^Ta_i+\phi_i^{(2)}(x)^Tb_i. \qquad(1)
\]

三维时，\(a_i\in\mathbb R^3\)、\(b_i\in\mathbb R^6\)，未减均值的基按如下顺序排列：

\[
B_i^{(1)}=\begin{bmatrix}\xi_x&\xi_y&\xi_z\end{bmatrix}^T,
\quad B_i^{(2)}=\begin{bmatrix}\xi_x^2/2&\xi_x\xi_y&\xi_x\xi_z&\xi_y^2/2&\xi_y\xi_z&\xi_z^2/2\end{bmatrix}^T. \qquad(2)
\]

分别定义 \(\mu_i^{(r)}=V_i^{-1}\int_{\Omega_i}B_i^{(r)}dV\)，令
\(\phi_i^{(r)}=B_i^{(r)}-\mu_i^{(r)}\)。不论系数怎样变化，都有
\(V_i^{-1}\int_{\Omega_i}P_i\,dV=\bar u_i\)。因此式（1）自动保持每个对偶体的已知均值。

节点处物理梯度 \(g_i\)、Hessian \(H_i\) 与归一化系数的关系为

\[
g_i=L_i^{-1}a_i,\qquad
H_i=L_i^{-1}\mathcal H(b_i)L_i^{-1},\qquad
\mathcal H(b)=\begin{bmatrix}b_1&b_2&b_3\\b_2&b_4&b_5\\b_3&b_5&b_6\end{bmatrix}. \qquad(3)
\]

三维``只迭代一次项''就是由每节点 \(9n_v\) 个系数未知量降到 \(3n_v\)
个；二维则由 \(5n_v\) 降到 \(2n_v\)。这里 \(n_v\)
为守恒变量数。下文以标量写公式，向量守恒量对同一几何算子并列求解。

\hypertarget{ux4eceux4f20ux7edfux53d8ux5206ux6cdbux51fdux5f97ux5230ux5206ux5757ux7cfbux7edf}{%
\section{3.
从传统变分泛函得到分块系统}\label{ux4eceux4f20ux7edfux53d8ux5206ux6cdbux51fdux5f97ux5230ux5206ux5757ux7cfbux7edf}}

王乾论文式（4-5）的 IJI 是界面上值及各阶法向导数跳跃平方的积分，第 4.2
节式（4-23）用 \(w\) 同时调整值和一阶导数项。当前 NCFV
代码采用边中点的值、完整梯度与 Hessian
跳跃，是受该思想启发的一种离散泛函；它与原论文的界面积分 IJI
有区别。以下消元推导对二者均适用，只需使用各自实际的残差行。

对于当前代码，原始边 \(e=(i,j)\) 长度为 \(\ell_e\)，中点为 \(m_e\)，令
\(c_i=(a_i,b_i)^T\)，定义

\[
D_i^e=\begin{bmatrix}
\sqrt w\,\phi_i(m_e)^T\\
\sqrt w\,(\ell_e/2)\nabla\phi_i(m_e)^T\\
(\ell_e^2/4)\operatorname{vec}(\nabla^2\phi_i)^T
\end{bmatrix},\qquad
r_e=D_i^e c_i-D_j^e c_j+\sqrt w(\bar u_i-\bar u_j)e_0. \qquad(4)
\]

其中最后一个矩阵按 \(d^2\) 个 Hessian
分量排列，混合导数重复计入，恰好得到 Frobenius
范数。每张界面只计一次，写

\[
J(c)=\tfrac12\sum_e\|r_e\|_2^2
=\tfrac12 c^TAc-f^Tc+\text{constant},\qquad Ac=f. \qquad(5)
\]

当前 \texttt{BuildVariationalOperators} 中没有额外的 \(A_e/\ell_e\)
面积权；本说明按其权重保持一致。若选择论文原式，可把积分点、法向导数和积分权重装入更多残差行，后续分块代数不变。

把全场系数按``一次项在前，二次项在后''重排：

\[
\begin{bmatrix}A_{aa}&A_{ab}\\A_{ba}&A_{bb}\end{bmatrix}
\begin{bmatrix}a\\b\end{bmatrix}
=\begin{bmatrix}f_a\\f_b\end{bmatrix},\qquad A_{ba}=A_{ab}^T. \qquad(6)
\]

三维 \(a\) 有 \(3N\) 个分量，\(b\) 有 \(6N\) 个分量。\(A_{bb}\)
是\textbf{全场二次项之间的耦合矩阵}，包含相邻单元/节点的二次项交叉块；它通常不是由独立的
\(6\times6\) 块组成。

从平方残差形式可知 \(A\) 对称半正定。以下涉及逆和普通 PCG 的论证均要求
\(A\)
在适用边界、周期约束及网格秩条件下正定。该条件应通过零空间分析和离散秩检查保证；仅检查每个局部对角块可逆不够。

对于当前二次中点泛函，一个充分条件的证明是：零能量使每条边两侧的
Hessian、边中点梯度及边中点值都相等；沿连通图传播，各零均值扰动多项式成为同一个全局二次多项式。若所有对偶体均值泛函对该多项式空间满列秩，就只能得到零多项式，故
\(A\) 正定。普通三维非周期域的完整二次空间维数为
10，需要几何独立的均值约束；周期情形要在一致展开和周期相容空间中检查，不能只数节点个数。

\hypertarget{ux4e3aux4ec0ux4e48ux4e0dux80fdux76f4ux63a5ux53eaux66f4ux65b0ux524dux4e09ux884c}{%
\section{4.
为什么不能直接只更新前三行}\label{ux4e3aux4ec0ux4e48ux4e0dux80fdux76f4ux63a5ux53eaux66f4ux65b0ux524dux4e09ux884c}}

\hypertarget{ux5220ux9664ux4e8cux6b21ux9879ux4f1aux6539ux53d8ux76eeux6807ux65b9ux7a0b}{%
\subsection{4.1
删除二次项会改变目标方程}\label{ux5220ux9664ux4e8cux6b21ux9879ux4f1aux6539ux53d8ux76eeux6807ux65b9ux7a0b}}

若直接置 \(b=0\)，式（6）被改成 \(A_{aa}a=f_a\)。原方程的一次项实际满足
\(A_{aa}a=f_a-A_{ab}b\)；被删掉的是二次项对梯度的反馈。原本二次精确的平均值关系一般也不再精确。

例如一维 \(u(x)=x^2/2\)，两个宽度均为 \(h\) 的控制体中心分别为
\(0,h\)，其均值为 \(h^2/24\) 与
\(h^2/2+h^2/24\)。若在边界处仅以这两个均值作线性斜率，则

\[
g_h=\frac{\bar u_1-\bar u_0}{h}=\frac h2,
\qquad u'(0)=0. \qquad(7)
\]

梯度误差为
\(O(h)\)。对称模板可能偶然消掉该误差，但一般非均匀、非结构网格没有这种保证。

\hypertarget{ux5355ux8282ux70b9ux5c40ux90e8-schur-ux6d88ux5143ux4ecdux9700ux8981ux90bbux70b9ux4e8cux6b21ux9879}{%
\subsection{4.2 单节点局部 Schur
消元仍需要邻点二次项}\label{ux5355ux8282ux70b9ux5c40ux90e8-schur-ux6d88ux5143ux4ecdux9700ux8981ux90bbux70b9ux4e8cux6b21ux9879}}

传统 block 扫描在节点 \(i\) 上解 \(K_i c_i=r_i\)，其中 \(r_i\)
含邻点当前系数。把局部块划分后可写

\[
S_i=K_{aa,i}-K_{ab,i}K_{bb,i}^{-1}K_{ba,i},\qquad
a_i=S_i^{-1}(r_{a,i}-K_{ab,i}K_{bb,i}^{-1}r_{b,i}). \qquad(8)
\]

但 \(r_{a,i}\)、\(r_{b,i}\) 仍包含邻点
\(b_j\)。为了让下一次扫描使用它们，还须恢复

\[
b_i=K_{bb,i}^{-1}(r_{b,i}-K_{ba,i}a_i). \qquad(9)
\]

因此局部块消元只改变局部求解组织方式，并未把全局迭代未知量降为一次项。同样，截取局部逆矩阵的前三行不能消去右端中的邻点二次系数。这与独立
LS
问题``完整二次矩阵求伪逆后仅输出梯度行''的情形不同：变分右端自身含待求的邻点多项式系数。

\hypertarget{ux8defux7ebf-aux4e25ux683cux7b49ux4ef7ux7684ux5168ux5c40-schur-ux6d88ux5143}{%
\section{5. 路线 A：严格等价的全局 Schur
消元}\label{ux8defux7ebf-aux4e25ux683cux7b49ux4ef7ux7684ux5168ux5c40-schur-ux6d88ux5143}}

由式（6）的第二行解出

\[
b_{\mathrm{opt}}(a)=A_{bb}^{-1}(f_b-A_{ba}a). \qquad(10)
\]

代入第一行，得到一次项系统

\[
Sa=\widetilde f_a,\qquad
S=A_{aa}-A_{ab}A_{bb}^{-1}A_{ba},\qquad
\widetilde f_a=f_a-A_{ab}A_{bb}^{-1}f_b. \qquad(11)
\]

从泛函看，\(b_{\mathrm{opt}}(a)\) 是固定 \(a\) 时使 \(J(a,b)\)
最小的二次项，故

\[
J_{\mathrm{red}}(a)=\min_b J(a,b)
=\tfrac12a^TSa-\widetilde f_a^Ta+\text{constant}. \qquad(12)
\]

式（11）求出的 \(a\)
与完整系统\textbf{充分收敛后的}一次项一致。它不保证与当前代码固定三轮扫描等未收敛结果相同；比较时应控制同一线性残差，而非固定相同扫描次数。

\hypertarget{ux552fux4e00ux6027ux548cux5bf9ux79f0ux6b63ux5b9aux6027}{%
\subsection{5.1
唯一性和对称正定性}\label{ux552fux4e00ux6027ux548cux5bf9ux79f0ux6b63ux5b9aux6027}}

若 \(A\succ0\)，其主子块 \(A_{bb}\succ0\)。对任意非零 \(p\)，令
\(q=(p,-A_{bb}^{-1}A_{ba}p)^T\)，则

\[
p^TSp=q^TAq>0. \qquad(13)
\]

因此 \(S\) 对称正定，一次项解唯一。这是严格消元的基本保证。若原 \(A\)
有零空间，应先处理边界约束或适用的商空间，不能直接调用此结论。

\hypertarget{ux4e0dux663eux5f0fux5f62ux6210ux7a20ux5bc6-schur-ux77e9ux9635}{%
\subsection{5.2 不显式形成稠密 Schur
矩阵}\label{ux4e0dux663eux5f0fux5f62ux6210ux7a20ux5bc6-schur-ux77e9ux9635}}

\(A_{bb}^{-1}\) 通常具有非局部作用，直接形成 \(S\)
会产生填充。一次矩阵自由乘法按下面三步执行：

\begin{Shaded}
\begin{Highlighting}[]
\NormalTok{ApplySchur(v):}
\NormalTok{    t = A\_ba * v}
\NormalTok{    solve A\_bb * z = t}
\NormalTok{    return A\_aa * v {-} A\_ab * z}
\end{Highlighting}
\end{Shaded}

每次重构的右端先解 \(A_{bb}z_f=f_b\)，再取
\(\widetilde f_a=f_a-A_{ab}z_f\)。几何不变时可以复用 \(A_{bb}\)
分解或预条件器。高效通量阶段只输出
\(g_i=L_i^{-1}a_i\)，无须为了通量再恢复最终的 \(b\)；但 Schur
算子内部仍有 \(6N\) 维临时量和线性求解。

在小网格上，可对 \(A_{bb}\) 用稀疏 Cholesky/LDLT
建立可信的等价基准。大网格可采用分块预条件或内部迭代。PETSc
的分块求解文档给出了这类 Schur
算子及内部子求解器的组织方式；这里引用其结构思想，并不要求项目引入
PETSc。{[}参考资料 3{]}

\hypertarget{ux5185ux90e8ux8fd1ux4f3cux6c42ux89e3ux5fc5ux987bux8ba1ux5165ux8befux5dee}{%
\subsection{5.3
内部近似求解必须计入误差}\label{ux5185ux90e8ux8fd1ux4f3cux6c42ux89e3ux5fc5ux987bux8ba1ux5165ux8befux5dee}}

设内部求解得到 \(\widetilde z\)，其残差
\(\epsilon=t-A_{bb}\widetilde z\)，则近似 Schur 作用误差为

\[
\widetilde S v-Sv=A_{ab}A_{bb}^{-1}\epsilon. \qquad(14)
\]

若每次内部求解用不同停止轮数，外层看到的作用未必是固定线性、对称算子；不能不加控制地套普通
PCG。FGMRES
允许变化的预条件器，也不会自动修复任意变化的主算子误差。可行策略是精确/足够准确地作用
\(A_{bb}^{-1}\)，或采用有误差控制的嵌套 Krylov 方法，并最终恢复
\(b_{\mathrm{opt}}(a)\) 核对完整系统残差。

直接把 \(A_{bb}^{-1}\) 换成局部块逆，所得方程是近似 Schur
方程，一般不再与原变分方法等价，正定性也需要另证。局部式（8）可作为外层预条件器的候选块，但它通常不是全局
\(S\) 的真实对角块。

\hypertarget{ux8defux7ebf-bux7528ux68afux5ea6ux5deeux6784ux9020-hessian-ux95edux5408}{%
\section{6. 路线 B：用梯度差构造 Hessian
闭合}\label{ux8defux7ebf-bux7528ux68afux5ea6ux5deeux6784ux9020-hessian-ux95edux5408}}

若目标是运行时只保留稀疏的一次项迭代，可让二次项成为一次项场的线性函数。这里给出一种可明确验证二次场精确性的构造，作为新的算法方案。

\hypertarget{ux4e8cux6b21ux51fdux6570ux7684ux68afux5ea6ux5deeux6052ux7b49ux5f0f}{%
\subsection{6.1
二次函数的梯度差恒等式}\label{ux4e8cux6b21ux51fdux6570ux7684ux68afux5ea6ux5deeux6052ux7b49ux5f0f}}

对任意二次多项式，其 Hessian 常数，故

\[
g_j-g_i=H_i(x_j-x_i). \qquad(15)
\]

为了与现有归一化基一致，定义 \(s_{ij}=L_i^{-1}(x_j-x_i)\)。由式（3）可得

\[
L_iL_j^{-1}a_j-a_i=\mathcal H(b_i)s_{ij}=B(s_{ij})b_i, \qquad(16)
\]

其中三维矩阵为

\[
B(s)=\begin{bmatrix}
s_x&s_y&s_z&0&0&0\\
0&s_x&0&s_y&s_z&0\\
0&0&s_x&0&s_y&s_z
\end{bmatrix}. \qquad(17)
\]

使用不同节点的归一化系数时，式（16）的 \(L_iL_j^{-1}\) 不能省略；直接把
\(a_j-a_i\) 当物理梯度差，在拉伸或尺度不一致的网格上会出错。

\hypertarget{ux81eaux7136ux90bbux57dfux4e0eux79e9ux6761ux4ef6}{%
\subsection{6.2
自然邻域与秩条件}\label{ux81eaux7136ux90bbux57dfux4e0eux79e9ux6761ux4ef6}}

初始设计取 \(\mathcal N_i\) 为所有与 \(i\)
直接由原始边相连的节点，不固定点数。将 \(B(s_{ij})\) 逐块叠成
\(B_i\in\mathbb R^{3|\mathcal N_i|\times6}\)，右端叠成 \(d_i(a)\)。设
\(\operatorname{rank}B_i=6\)，预计算其 SVD 伪逆，定义

\[
b_i=B_i^\dagger d_i(a)=T_i a,
\qquad b=Ta. \qquad(18)
\]

具体存储时，把 \(B_i^\dagger\) 按邻点拆为 \(6\times3\) 子块
\(Q_{ij}\)，则可直接预计算

\[
T_{ij}=Q_{ij}L_iL_j^{-1}\ (j\in\mathcal N_i),\qquad
T_{ii}=-\sum_{j\in\mathcal N_i}Q_{ij}, \qquad(18b)
\]

其他块为零。运行时只需邻点矩阵乘加，不再做 SVD；二维同理使用
\(3\times2\) 子块。

本方案初版使用等权拟合，保留 \(w\) 为原泛函中值项与导数项的相对权重；SVD
容差是数值秩判据。若邻点位移张成三维空间，则 \(B_i h=0\) 推出
\(\mathcal H(h)\) 在三个独立方向上作用为零，从而 \(h=0\)，所以满列秩。

若几何退化、邻域不足或周期最短位移造成秩亏，必须明确报告。可将邻域扩展为所属原始单元全部顶点或下一拓扑环，再重新查秩；这种扩展应由自然邻接和秩条件决定，不截取任意固定数量的点。后续
halo 范围须随实际邻域更新。

对精确二次函数，式（16）每行成立，满秩时式（18）精确恢复
\(b_i\)。对一般光滑函数，\(b_i\)
是梯度差给出的局部近似。它不是式（10）的全局最小化二次项，二者应分清。

\hypertarget{ux4ee3ux56deux6cdbux51fdux540eux5fc5ux987bux91cdux65b0ux6c42ux5bfc}{%
\section{7.
代回泛函后必须重新求导}\label{ux4ee3ux56deux6cdbux51fdux540eux5fc5ux987bux91cdux65b0ux6c42ux5bfc}}

定义满列秩嵌入 \(E=(I,T)^T\)，于是 \(c=Ea\)。新的受约束泛函为

\[
\widehat J(a)=J(Ea),\qquad
\widehat A a=\widehat f,\qquad
\widehat A=E^TAE,\quad \widehat f=E^Tf. \qquad(19)
\]

展开得到

\[
\widehat A=A_{aa}+A_{ab}T+T^TA_{ba}+T^TA_{bb}T,
\qquad \widehat f=f_a+T^Tf_b. \qquad(20)
\]

这是路线 B 的核心方程。若只在式（6）的第一行代 \(b=Ta\)，会得到
\((A_{aa}+A_{ab}T)a=f_a\)，漏掉 \(T^T\) 反馈；该式一般非对称，也不是
\(\widehat J\) 的驻值方程。

\hypertarget{ux7a00ux758fux6b8bux5deeux884cux7684ux76f4ux63a5ux6784ux9020}{%
\subsection{7.1
稀疏残差行的直接构造}\label{ux7a00ux758fux6b8bux5deeux884cux7684ux76f4ux63a5ux6784ux9020}}

将式（4）的端点算子按一次、二次列分为 \(D_i^e=[D_{ia}^e,D_{ib}^e]\)。用
\(T_{ik}\) 表示 \(b_i\) 对节点 \(k\) 的 \(a_k\) 的线性块，则

\[
\widehat D_{ek}=D_{ia}^e\mathbf1_{k=i}-D_{ja}^e\mathbf1_{k=j}
+D_{ib}^eT_{ik}-D_{jb}^eT_{jk}. \qquad(21)
\]

仅当 \(k\) 位于两个端点或其 Hessian 恢复邻域时该块非零。令
\(q_e=\sqrt w(\bar u_i-\bar u_j)e_0\)，可直接写

\[
\widehat J(a)=\tfrac12\sum_e\left\|q_e+\sum_k\widehat D_{ek}a_k\right\|_2^2. \qquad(22)
\]

相应的节点方程为

\[
\sum_l\left(\sum_e\widehat D_{ek}^T\widehat D_{el}\right)a_l
=-\sum_e\widehat D_{ek}^Tq_e. \qquad(23)
\]

这允许只保存每条残差的 \(3\)
列系数块和支撑列表，按``先算残差，再转置回加''的方式实现矩阵自由乘法，避免显式构造完整块矩阵。

\hypertarget{ux6b63ux5b9aux6027ux7cbeux786eux6027ux4e0eux539fux65b9ux6cd5ux7684ux5deeux522b}{%
\subsection{7.2
正定性、精确性与原方法的差别}\label{ux6b63ux5b9aux6027ux7cbeux786eux6027ux4e0eux539fux65b9ux6cd5ux7684ux5deeux522b}}

若 \(A\succ0\)，由于 \(E\) 的上半部分为单位阵，对任意 \(a\ne0\) 有

\[
a^T\widehat Aa=(Ea)^TA(Ea)>0. \qquad(24)
\]

所以新系统对称正定且有唯一解。若原系统仅半正定，准确条件是
\(\ker A\cap\operatorname{range}E=\{0\}\)。局部 Hessian
拟合满秩本身不足以证明全局正定。

若原残差对全局二次场精确、所有 \(B_i\)
满秩、均值与周期展开一致，则精确二次场满足 \(b=Ta\)
且全部残差为零，因而新方法仍能精确再现该二次场。但二次再现只是三阶空间精度的必要基础；还需网格族上的稳定性、边界一致性、充分收敛和通量/时间误差控制。

特别地，每张固定网格上正定不等于加密序列上归一化逆算子一致有界。若闭合邻域越来越扁平或矩阵条件数失控，误差常数仍可能随加密增长；需要以几何秩、缩放后的条件数和实际加密结果共同判断。

设完整系统极小点为 \(c_{\mathrm{opt}}\)，新系统极小点为
\(E\widehat a\)。二次型的能量恒等式为

\[
J(E\widehat a)-J(c_{\mathrm{opt}})
=\tfrac12(E\widehat a-c_{\mathrm{opt}})^TA(E\widehat a-c_{\mathrm{opt}})\ge0. \qquad(25)
\]

这表明新方法是在受约束空间中的最优解，一般不会与传统变分梯度完全相同。它保留的是可证明的二次场一致性和变分结构；对一般流场的误差大小与鲁棒性需要实测。

\hypertarget{ux7ebfux6027ux6c42ux89e3ux6570ux636eux548c-mpi-ux901aux4fe1}{%
\section{8. 线性求解、数据和 MPI
通信}\label{ux7ebfux6027ux6c42ux89e3ux6570ux636eux548c-mpi-ux901aux4fe1}}

\hypertarget{ux5bf9ux8defux7ebf-b-ux4f7fux7528ux4e09ux4e58ux4e09ux5757ux9884ux6761ux4ef6-pcg}{%
\subsection{8.1 对路线 B 使用三乘三块预条件
PCG}\label{ux5bf9ux8defux7ebf-b-ux4f7fux7528ux4e09ux4e58ux4e09ux5757ux9884ux6761ux4ef6-pcg}}

由式（23）取真实对角块

\[
K_k=\sum_e\widehat D_{ek}^T\widehat D_{ek},\qquad
M^{-1}=\operatorname{diag}(K_k^{-1}). \qquad(26)
\]

每个 \(K_k\) 为 \(3\times3\)（二维为 \(2\times2\)），初始化时 LDLT
分解并检查正性。所有守恒变量复用几何块。PCG 依次使用
\(r=f-Ax\)、\(z=M^{-1}r\)、\(p=z\)，每轮

\[
\alpha=\frac{r^Tz}{p^TAp},\quad x\leftarrow x+\alpha p,
\quad r_{\mathrm{new}}=r-\alpha Ap,\quad z_{\mathrm{new}}=M^{-1}r_{\mathrm{new}}. \qquad(27)
\]

\[
\beta=\frac{r_{\mathrm{new}}^Tz_{\mathrm{new}}}{r^Tz},\qquad
p\leftarrow z_{\mathrm{new}}+\beta p. \qquad(28)
\]

这里 \(A,f\) 指
\(\widehat A,\widehat f\)。建议同时记录相对预条件残差、真实残差及最大迭代数；上限耗尽应给出失败状态。设计初期可取相对于初始残差和右端预条件范数中较大者的
\(10^{-10}\)
阈值，再用误差预算确认，无须把该数值当作所有网格的理论常量。

\hypertarget{ux4e24ux79cdux7b49ux4ef7ux7684ux7b97ux5b50ux5b9eux73b0}{%
\subsection{8.2
两种等价的算子实现}\label{ux4e24ux79cdux7b49ux4ef7ux7684ux7b97ux5b50ux5b9eux73b0}}

\textbf{预计算残差块。} 储存式（21）的
\(\widehat D_{ek}\)；每轮同步所有支撑上的 \(a_k\)，逐残差计算
\(s_e=\sum_k\widehat D_{ek}a_k\)，再对本 rank 的 owned 行累加
\(\widehat D_{ek}^Ts_e\)。每个 owned
节点必须拥有所有对其有贡献的残差行副本。只保留本 rank
拥有的原始边会漏掉非端点支撑节点的方程贡献。

\textbf{分阶段复合算子。} 不展开 \(\widehat D\)，按下述过程作用
\(E^TAE\)：

\begin{Shaded}
\begin{Highlighting}[]
\NormalTok{pull ghost values of a}
\NormalTok{b\_owned = T * a}
\NormalTok{pull ghost values of b}
\NormalTok{(y\_a, y\_b) = A * (a, b)}
\NormalTok{z = y\_a}
\NormalTok{scatter\_add T\^{}T * y\_b to owners of a}
\NormalTok{return z}
\end{Highlighting}
\end{Shaded}

最后一步是转置累加，远端贡献要回加到拥有者。只做普通 ghost Pull
无法代替它，否则实际实现会漏项并破坏对称性。预计算残差块可以改用较宽的
Pull 避免运行时转置回加；两种路线应由内存与通信性能决定。

\hypertarget{ux65b0ux65b9ux6cd5ux4e0dux80fdux76f4ux63a5ux6cbfux7528ux4e24ux73afux5373ux53efux7684ux7ed3ux8bba}{%
\subsection{8.3
新方法不能直接沿用``两环即可''的结论}\label{ux65b0ux65b9ux6cd5ux4e0dux80fdux76f4ux63a5ux6cbfux7528ux4e24ux73afux5373ux53efux7684ux7ed3ux8bba}}

若原 \(A\) 耦合原始边一跳、\(T\) 使用一环边邻点，则 \(T^TA_{bb}T\)
在节点图上最多连接三跳：\(k\to i\to j\to l\)。因此显式存储凝聚矩阵行时可能需要三跳格点依赖。若
\(T\) 使用更宽的单元星邻域，实际范围还会变化。

分阶段实现可通过多次较窄的通信完成相同算子，无须一次持有所有三跳未知量。几何单元
ghost 与求解变量 halo 仍应分别按所需几何和非零支撑收集。路线 A 的精确
Schur 作用是非局部的，由内部全局求解传播信息，也不等同于固定两环通信。

\hypertarget{ux4e3aux4ec0ux4e48ux9ad8ux6548ux79efux5206ux53efux4ee5ux53eaux4f7fux7528ux68afux5ea6}{%
\section{9.
为什么高效积分可以只使用梯度}\label{ux4e3aux4ec0ux4e48ux9ad8ux6548ux79efux5206ux53efux4ee5ux53eaux4f7fux7528ux68afux5ea6}}

\hypertarget{ux5728ux5faeux5355ux7eafux5f62ux4e0aux6d88ux53bbux79efux5206ux516cux5f0fux4e2dux7684-hessian}{%
\subsection{9.1 在微单纯形上消去积分公式中的
Hessian}\label{ux5728ux5faeux5355ux7eafux5f62ux4e0aux6d88ux53bbux79efux5206ux516cux5f0fux4e2dux7684-hessian}}

设微单纯形 \(K\) 有 \(m\) 个顶点 \(y_a\)（三角形 \(m=3\)，四面体
\(m=4\)），锚点为 \(x_i\)，\(d_a=y_a-x_i\)，\(D=\sum_a d_a\)。对二次函数

\[
u(x)=u_i+g_i\cdot(x-x_i)+\tfrac12(x-x_i)^TH(x-x_i), \qquad(29)
\]

单纯形矩公式给出

\[
\int_Ku=|K|u_i+\frac{|K|}{m}D\cdot g_i
+\frac{|K|}{2m(m+1)}\left[D^THD+\sum_a d_a^THd_a\right]. \qquad(30)
\]

二次函数满足 \(Hd_a=g(y_a)-g_i\)。代入并合并同类项：

\[
\int_Ku=|K|u_i+\frac{|K|}{2m}D\cdot g_i
+\frac{|K|}{2m(m+1)}\sum_a(D+d_a)\cdot g(y_a). \qquad(31)
\]

若构造点具有仿射支撑
\(y_a=\sum_k C_{ak}x_k\)、\(\sum_kC_{ak}=1\)，则对二次函数的线性梯度场
\(g(y_a)=\sum_kC_{ak}g_k\)。可把式（31）写成

\[
\int_Ku=|K|u_i+\sum_k\ell_{ik}^K\cdot g_k,\quad
\ell_{ik}^K=\frac{|K|}{2m}D\mathbf1_{k=i}
+\frac{|K|}{2m(m+1)}\sum_a C_{ak}(D+d_a). \qquad(32)
\]

几何固定时预计算
\(\ell_{ik}^K\)，每个时间阶段只做节点梯度的线性组合。这一步消去的是\textbf{积分中的
Hessian 表达}，不代替第 5---7 节对重构方程的消元或闭合。

\hypertarget{ux4eceux4f53ux5747ux503cux6062ux590dux8282ux70b9ux70b9ux503c}{%
\subsection{9.2
从体均值恢复节点点值}\label{ux4eceux4f53ux5747ux503cux6062ux590dux8282ux70b9ux70b9ux503c}}

对 \(\Omega_i\) 内的微体求和，定义
\(\omega^V_{ik}=V_i^{-1}\sum_{K\subset\Omega_i}\ell_{ik}^K\)，于是

\[
u_i^E=\bar u_i-\sum_k\omega^V_{ik}\cdot g_k. \qquad(33)
\]

当前 \texttt{pointRecoveryStencil} 已存储非零
\(\omega^V_{ik}\)，\texttt{RecoverPointValues} 可直接使用路线 A 或 B
输出的梯度。因此不需要为高效通量恢复最终的 \(b_i\)。

原局部多项式点值则为
\(P_i(x_i)=\bar u_i-\mu_i^{(1)T}a_i-\mu_i^{(2)T}b_i\)。在二次场精确条件下它与式（33）相同；对一般数据二者不是严格相同的离散映射。所以即使
Schur
路线给出完全相同的梯度，``传统积分''和``高效积分''的最终流场误差仍可能不同。

在形状规则的网格族上，若梯度误差为
\(O(h^2)\)、\(\omega^V_{ik}=O(h)\)、支撑与权重受控，则式（33）的点值误差为
\(O(h^3)\)，再加上光滑场二次积分的 \(O(h^3)\)
平均误差。该局部结论还不足以单独证明完整有限体积半离散格式的全局三阶收敛。

\hypertarget{ux5de6ux53f3ux5b8fux9762ux5747ux503c}{%
\subsection{9.3
左右宏面均值}\label{ux5de6ux53f3ux5b8fux9762ux5747ux503c}}

对原始边 \(e=(i,j)\) 的对偶宏面，记面积为 \(A_e\)、已有两侧梯度权重为
\(W_{ek}^{L/R}\)。无限制时

\[
\bar u_{e,L}=u_i^E+\frac1{A_e}\sum_kW_{ek}^L\cdot g_k,
\qquad
\bar u_{e,R}=u_j^E+\frac1{A_e}\sum_kW_{ek}^R\cdot g_k. \qquad(34)
\]

将式（33）代入，记左右面均值的差为
\(\delta_e\)，整个跳跃也可仅用梯度写成

\[
\delta_e=\bar u_i-\bar u_j
+\sum_k\left[-\omega^V_{ik}+\omega^V_{jk}
+\frac{W_{ek}^{L}-W_{ek}^{R}}{A_e}\right]\cdot g_k. \qquad(35)
\]

未出现的权重视为零。式（35）解释梯度与高效面公式的完整耦合；本文路线 A/B
的目标仍是传统泛函及其受约束版本，并未额外把式（35）改造成另一套面均值目标。

\hypertarget{ux7269ux7406ux901aux91cfux5b88ux6052ux548c-ssprk3}{%
\section{10. 物理通量、守恒和
SSPRK3}\label{ux7269ux7406ux901aux91cfux5b88ux6052ux548c-ssprk3}}

恢复保守点值 \(U_i^E\) 与梯度 \(g_{i\alpha}=\partial_\alpha U_i\)
后，物理通量梯度通过链式法则计算：

\[
\partial_\alpha F_\beta(U_i^E)
=\frac{\partial F_\beta}{\partial U}(U_i^E)g_{i\alpha}. \qquad(36)
\]

令
\(\boldsymbol A_e=\int_{\Gamma_e}n\,dS\)，则已有几何权重给出两侧物理通量积分

\[
I^F_{e,s}=F(U_{a_s}^E)\cdot\boldsymbol A_e
+\sum_{k,\alpha,\beta}W_{ek}^{s,\alpha\beta}
\partial_\alpha F_\beta(U_k^E),\qquad a_L=i,\ a_R=j. \qquad(37)
\]

对于一般非线性 Euler
通量，这是一种高阶积分近似；标量二次场上的式（32）精确性不能直接解释为任意非线性通量严格精确。当前高效算子用宏面代表单位法向
\(\widehat n_e=\boldsymbol A_e/\|\boldsymbol A_e\|\)。定义面均值处的中心通量
\(C_e\) 与数值耗散修正 \(\mathcal D_e\)：

\[
C_e=\tfrac12\left[F(\bar U_{e,L})+F(\bar U_{e,R})\right]\cdot\widehat n_e. \qquad(38a)
\]

\[
\mathcal D_e=A_e\left[\widehat F(\bar U_{e,L},\bar U_{e,R},\widehat n_e)-C_e\right]. \qquad(38b)
\]

最终宏面积分为

\[
\widehat I_e=\tfrac12(I^F_{e,L}+I^F_{e,R})+\mathcal D_e. \qquad(38c)
\]

平面宏面时
\(\|\boldsymbol A_e\|=A_e\)；一般折面宏面不能把两者混用。式（38）的几何选择与现有
\texttt{EvaluateOwnedEdgeFlux}
保持一致。它每张宏面只调用一次黎曼解算器，以平均物理积分作中心项，再加入面均值上的数值耗散。对两个节点分别累加
\(+\widehat I_e\) 与 \(-\widehat I_e\)，维持全局守恒。

设
\(L_i(U)=-V_i^{-1}\sum_e\sigma_{ie}\widehat I_e\)（另加边界、黏性或源项），SSPRK3
为

\[U^{(1)}=U^n+\Delta tL(U^n).\qquad(39a)\]

\[U^{(2)}=\tfrac34U^n+\tfrac14[U^{(1)}+\Delta tL(U^{(1)})].\qquad(39b)\]

\[U^{n+1}=\tfrac13U^n+\tfrac23[U^{(2)}+\Delta tL(U^{(2)})].\qquad(39c)\]

每次阶段右端均重新用当前均值求一次项场，并完成梯度、点值和必要通量数据的
halo
同步。限制器可沿用高效模式的每侧单一系数；但限制器或非物理状态回退激活时，应分别讨论线性变分最优性和实际通量状态，不能继续声称整个流程是同一个线性极小化问题。

\hypertarget{ux53efux5b9eux65bdux7684ux7a0bux5e8fux6b65ux9aa4}{%
\section{11.
可实施的程序步骤}\label{ux53efux5b9eux65bdux7684ux7a0bux5e8fux6b65ux9aa4}}

\hypertarget{ux8defux7ebf-a-ux7684ux4e25ux683cux7b49ux4ef7ux57faux51c6}{%
\subsection{11.1 路线 A
的严格等价基准}\label{ux8defux7ebf-a-ux7684ux4e25ux683cux7b49ux4ef7ux57faux51c6}}

\begin{Shaded}
\begin{Highlighting}[]
\NormalTok{初始化：}
\NormalTok{  从当前变分残差行建立未经左预条件的 A\_aa,A\_ab,A\_ba,A\_bb}
\NormalTok{  为 A\_bb 建立分解/受控子求解器；为 Schur 系统建预条件器}

\NormalTok{每次重构：}
\NormalTok{  同步当前控制体均值；装配 f\_a,f\_b}
\NormalTok{  求 A\_bb*z\_f=f\_b；得到 f\_tilde=f\_a{-}A\_ab*z\_f}
\NormalTok{  在 ApplySchur 上求一次项 a}
\NormalTok{  必要时恢复 b 并核对完整残差 A*(a,b){-}f}
\NormalTok{  输出 g\_i=L\_i\^{}\{{-}1\}a\_i；进入现有高效积分流程}
\end{Highlighting}
\end{Shaded}

现有 \texttt{\_variationalOperators} 存的是局部 \(K_i^{-1}C_{ij}\) 与
\(K_i^{-1}\)
作用后的均值项，已做左块预条件。不能直接把这些存储块当成欧氏内积下对称的原始
\(A\) 来使用 PCG 或套式（13）；应保留原始残差行或重新构造对称矩阵。

\hypertarget{ux8defux7ebf-b-ux7684ux68afux5ea6ux95edux5408ux7248ux672c}{%
\subsection{11.2 路线 B
的梯度闭合版本}\label{ux8defux7ebf-b-ux7684ux68afux5ea6ux95edux5408ux7248ux672c}}

\begin{Shaded}
\begin{Highlighting}[]
\NormalTok{几何初始化：}
\NormalTok{  建立原变分残差行 D\_i\^{}e，并区分一次列和二次列}
\NormalTok{  对每个 owned 节点收集自然邻域，构建 B\_i 并做秩揭示 SVD}
\NormalTok{  预计算 T\_ik；选择残差块展开或分阶段算子实现}
\NormalTok{  收集完整矩阵行及所有 ghost 支撑}
\NormalTok{  预计算三乘三对角块 K\_i，并检查正性}

\NormalTok{每次 EvaluateRHS：}
\NormalTok{  同步控制体均值，形成各残差已知项 q\_e}
\NormalTok{  f\_hat\_k = {-}sum\_e D\_hat\_ek\^{}T*q\_e}
\NormalTok{  以上次梯度或完整二次 LS 的一次项为初值}
\NormalTok{  用 PCG 求 A\_hat*a=f\_hat；检查残差和最大迭代数}
\NormalTok{  g\_i=L\_i\^{}\{{-}1\}a\_i；同步梯度 halo}
\NormalTok{  利用 pointRecoveryStencil 恢复并同步点值}
\NormalTok{  按已有高效路径处理限制器、物理通量梯度和宏面通量}
\NormalTok{  守恒累加 RHS，完成本次 SSPRK3 阶段}
\end{Highlighting}
\end{Shaded}

为了保持二次精确性，LS 初值若采用现有二次
LS，应先建立完整二次问题，再取其一次项输出；初值不会改变充分收敛后的式（23）。此初值也可以直接取上次重构结果，算法本身不要求固定模板点数。

\hypertarget{ux4e0eux73b0ux6709ux6a21ux5757ux7684ux63a5ux53e3ux5bf9ux5e94}{%
\subsection{11.3
与现有模块的接口对应}\label{ux4e0eux73b0ux6709ux6a21ux5757ux7684ux63a5ux53e3ux5bf9ux5e94}}

\begin{longtable}[]{@{}
  >{\raggedright\arraybackslash}p{(\columnwidth - 2\tabcolsep) * \real{0.5000}}
  >{\raggedright\arraybackslash}p{(\columnwidth - 2\tabcolsep) * \real{0.5000}}@{}}
\toprule\noalign{}
\begin{minipage}[b]{\linewidth}\raggedright
模块
\end{minipage} & \begin{minipage}[b]{\linewidth}\raggedright
需要增加或复用的内容
\end{minipage} \\
\midrule\noalign{}
\endhead
\bottomrule\noalign{}
\endlastfoot
\texttt{NCFVReconstruction} &
新增降维算子、一次项场和求解流程；复用基、矩、参考长度 \\
\texttt{NCFVNodeHalo} & 按真实复合支撑收集依赖，或支持转置回加通信 \\
\texttt{NCFVSpatial::Reconstruct} & 接收输出梯度，复用梯度
Pull、点值恢复与后续阶段 \\
\texttt{NCFVDualGeometry} &
复用体积恢复及两侧面状态、物理通量积分权重 \\
\texttt{NCFVConfig} & 新增方法选择、残差容差、最大轮数、秩亏处理策略 \\
时间循环 & 继续使用当前 SSPRK3，每阶段重算重构与通量 \\
\end{longtable}

建议未来把三种重构清楚命名：现有 \texttt{Variational}；严格等价实验选项
\texttt{VariationalSchurGradient}；新闭合方法
\texttt{VariationalGradientClosure}。后两者只是\textbf{建议名称，当前
JSON 尚不接受这些选项}。\texttt{variationalWeight}
继续控制同一泛函中的值及一次导数权重，不能与 PCG 容差或 Hessian
拟合秩阈值混用。

为实际减少常驻系数内存，可给新模式增加只输出梯度的求解入口，使高效路径不再常驻分配每节点九行的
\texttt{\_coefficients}。路线 B
若采用分阶段作用，六行二次项是可复用的临时工作区；若采用式（21）的残差块展开，则可进一步省去该临时场。需检查高效限制器和点值恢复分支确实只读取梯度，避免接口仍强制分配完整二次场而抵消存储收益。

\hypertarget{ux7cbeux5ea6ux6027ux80fdux4e0eux5b9eux65bdux987aux5e8f}{%
\section{12.
精度、性能与实施顺序}\label{ux7cbeux5ea6ux6027ux80fdux4e0eux5b9eux65bdux987aux5e8f}}

一次项未知量减少，不自动等于运行时间或内存按同样比例减少。路线 A
仍有二次项内部解；路线 B 需要
\(T\)、更宽的残差支撑及相应通信。以下表格说明各路线真正减少的对象。

\begin{longtable}[]{@{}llll@{}}
\toprule\noalign{}
方案 & 三维每节点未知量 & 与原收敛重构等价 & 主要代价 \\
\midrule\noalign{}
\endhead
\bottomrule\noalign{}
\endlastfoot
完整变分后提取梯度 & 9 & 是 & 全系数扫描与通信 \\
每节点局部 Schur & 仍需 9 & 是 & 邻点二次项仍不可丢 \\
全局 Schur & 外层 3、内部 6 & 精确求解时是 & 全局子求解与填充 \\
梯度闭合 \(E^TAE\) & 3 & 一般不同 & 闭合算子和支撑扩展 \\
\end{longtable}

建议先用路线 A
建小规模代数基准，确认一次项与完整系统在相同残差标准下相符；再独立实现路线
B，验证二次一致性、MPI
对称性和守恒。若直接追求低存储且愿意评估新的离散方法，路线 B
更符合``仅以一次项为迭代未知量''的目标。

实施后的必要验证应包括：

\begin{enumerate}
\def\labelenumi{\arabic{enumi}.}
\tightlist
\item
  常数、一次和全部二次单项式；输入应为解析对偶体均值，检查
  \(a\)、梯度及点值，不只检查时间推进后的综合误差。二次多项式测试使用非周期域或一致的局部展开，不能强行让一般
  \(x^2\) 场满足周期约束。
\item
  路线 A
  与充分收敛完整系统比较一次项和完整残差，控制内部求解误差；不以固定三轮扫描为等价标准。
\item
  路线 B 检查 \(b=Ta\) 的二次精确性、随机向量的离散对称性
  \(x^T\widehat Ay-y^T\widehat Ax\)、非正曲率和秩亏诊断。
\item
  MPI 1/2/4/8 rank
  对比，覆盖支撑跨分区、周期展开和转置回加；核对结果不随分区改变。
\item
  固定足够小时间步做空间加密，再固定网格做时间加密，区分梯度、节点恢复、通量和时间误差。
\item
  再运行三棱柱四级等熵涡到
  \(t=1.6\)，报告真实残差、迭代次数、质量漂移、耗时和峰值内存。
\end{enumerate}

本工作区已有``完整变分迭代后输出一次导数 +
高效积分''的测试报告；其中的数据只能作为路线实施前的组合基线，不能当作新路线
A/B 的验证结果。

\hypertarget{ux53c2ux8003ux8d44ux6599ux4e0eux4ee3ux7801ux4f9dux636e}{%
\section{参考资料与代码依据}\label{ux53c2ux8003ux8d44ux6599ux4e0eux4ee3ux7801ux4f9dux636e}}

\begin{enumerate}
\def\labelenumi{\arabic{enumi}.}
\tightlist
\item
  王乾：《非结构网格紧致高精度有限体积方法》，用户提供的博士论文，第 4.1
  节的 IJI、稀疏系统及正定性讨论；第 4.2 节式（4-23）的权重 \(w\)。本文
  Schur 分块及梯度闭合为针对 NCFV 的推导。
\item
  当前工作区
  \texttt{src/NCFV/NCFVReconstruction.cpp}（变分重构、系数提取、点值恢复）与
  \texttt{src/NCFV/NCFVSpatial.cpp}（高效通量和调用顺序）。关键函数及接口见第
  1、3、9---11 节。核对日期：2026-09-28。
\item
  \href{https://petsc.org/release/manual/ksp/\#solving-block-matrices}{PETSc
  官方文档：Solving Block Matrices}，Schur
  算子、分块预条件与内部子求解器说明。本文未引入 PETSc 依赖。
\item
  项目内高效积分推导：\texttt{docs/guides/ncfv\_efficient\_differential\_porting\_zh.md}，第
  5---8 节。
\item
  当前组合基线：\texttt{docs/reports/ncfv\_efficient\_variational\_gradient\_t1p6\_dt\_20260928/report.md}。该报告仍使用完整二次变分扫描，未验证本文提出的新降维系统。
\end{enumerate}

\end{document}
